\chapter{Solicitação de registro de especialista}
\label{chap:pf-especialista}
\vfill

\section{O que é?}
\vfill

O registro de especialista atesta a experiência profissional em uma ou mais \href{https://site.cfp.org.br/servicos/titulo-de-especialista/}{áreas de especialidade} reconhecidas pelo Conselho Federal de Psicologia (CFP) e não constitui condição obrigatória para o exercício profissional.
\vfill

O registro de especialista do CFP não é o mesmo que o título acadêmico conferido a quem conclui uma pós-graduação \emph{lato sensu}:
\vfill

\begin{itemize}
	\item
	\textbf{especialista} é a/o profissional graduada/o que acumula notório saber teórico-prático em uma área específica, reconhecida por seu conselho de classe. Esse reconhecimento favorece a institucionalização da profissão na sociedade brasileira e ajuda a demarcar suas áreas de atuação;
	\item
	\textbf{especialização} é o aprimoramento profissional obtido por meio de curso de pós-graduação em uma área específica, que atenda às normativas e aos critérios de funcionamento determinados pelo MEC ou outro órgão com essa competência.
\end{itemize}
\vfill

Desse modo, há diferença entre ser psicóloga/o especialista e psicóloga/o com especialização: enquanto a/o primeira/o tem sua formação
e prática profissional reconhecidas/os por seu conselho regional, a/o segunda/o concluiu um curso de especialização.
\vfill

\section{Quem pode utilizar esse serviço?}

Qualquer psicóloga/o que possua o CRP ativo, com no mínimo 2 (anos) de inscrição, que esteja adimplente com as anuidades dos exercícios
anteriores e que não esteja cumprindo penalidade resultante de processo ético de multa, suspensão ou cassação.
\vfill

\section{Como solicitar o	serviço?}

Os pedidos são recebidos diretamente pelos \href{https://crpsp.org/sede/index}{\emph{e-mail}s institucionais das subsedes}.
\vfill\clearpage

\section{Que informações ou documentos são necessários?}

\begin{itemize}
	\item
	Certificado e histórico de curso de pós-graduação (reconhecido pelo MEC) em alguma das
	\href{https://site.cfp.org.br/servicos/titulo-de-especialista/}{áreas definidas pelo CFP}; ou
	\item
	documento que comprove a homologação do resultado final de concurso para registro de especialista promovido pelo CFP;
	\item
	comprovação de exercício profissional de no mínimo dois anos na	especialidade requerida (para cursos e provas concluídos a partir de 2
	de fevereiro de 2023);
	\item
	\href{http://www.crpsp.org.br/arquivos/SP-Formulario-de-Requerimento-CARPE.pdf}{requerimento de	registro de especialista} preenchido e assinado.
\end{itemize}

Mais informações estão disponíveis no \href{http://www.crpsp.org/pagina/view/51}{\emph{site} do CRP~SP}.
\vfill

\section{Quais são as etapas para a realização desse serviço?}

\begin{enumerate}
	\item
	A/o psicóloga/o solicita o registro de especialista via \emph{e-mail},
	anexando o
	\href{http://www.crpsp.org.br/arquivos/SP-Formulario-de-Requerimento-CARPE.pdf}{requerimento de	registro de especialista} e cópias (escaneadas) dos documentos	necessários para cada situação;
	\item
	os documentos são analisados pela Comissão de Análise para Concessão de Registro de Psicóloga/o Especialista (Carpe);
	\item
	a Plenária do CRP~vota pelo deferimento ou indeferimento do registro.
\end{enumerate}
\vfill

\section{Quanto custa?}

Gratuito. Será cobrada posteriormente taxa de emissão de nova CIP com o acréscimo de registro de especialista, caso o pedido seja deferido.
\vfill

\section{Quanto tempo leva?}

60 dias a partir da data de protocolo do documento, conforme \href{https://atosoficiais.com.br/cfp/resolucao-do-exercicio-profissional-n-23-2022-institui-condicoes-para-concessao-e-registro-de-psicologa-e-psicologo-especialistas-reconhece-as-especialidades-da-psicologia-e-revoga-as-resolucoes-cfp-n-13-de-14-de-setembro-de-2007-n-3-de-5-de-fevereiro-de-2016-n-18-de-5-de-setembro-de-2019}{Resolução CFP nº~23/2022}.

Caso a documentação esteja incompleta, a Carpe solicitará regularização, estendendo o período de análise final em 30 dias.

Em caso de indeferimento, é possível interpor recurso, devendo a/o psicóloga/o manifestar interesse em até 30 dias de ciência do
indeferimento. Não há prazo para análise do pedido de recurso pelo CFP.
\vfill

\section{Como ter acesso ao atual \emph{status} do serviço?}

O andamento do pedido pode ser consultado por \emph{e-mail}, telefone ou presencialmente. Para casos de recurso, as informações são
disponibilizadas pelo CFP a partir da identificação do protocolo SEI! vinculado ao pedido e comunicado à/ao psicóloga/o pelo CRP~SP.
\vfill\clearpage

\section{Legislação relacionada}

\begin{itemize}
	\item
	\href{https://site.cfp.org.br/wp-content/uploads/2008/08/Resolucao_CFP_nx_013-2007.pdf}{Resolução CFP nº~13/2007}, que institui a Consolidação das Resoluções relativas ao Título Profissional de Especialista em Psicologia e dispõe sobre normas e procedimentos para seu registro;
	\item
	\href{https://atosoficiais.com.br/cfp/resolucao-do-exercicio-profissional-n-23-2022-institui-condicoes-para-concessao-e-registro-de-psicologa-e-psicologo-especialistas-reconhece-as-especialidades-da-psicologia-e-revoga-as-resolucoes-cfp-n-13-de-14-de-setembro-de-2007-n-3-de-5-de-fevereiro-de-2016-n-18-de-5-de-setembro-de-2019}{Resolução CFP nº~23/2022}, que institui condições para concessão e registro de	psicóloga e psicólogo especialistas e reconhece as especialidades da Psicologia.
\end{itemize}

\section{Serviços correlatos}

\begin{itemize}
	\item
	\hyperref[chap:pf-cip]{Entrega da carteira de identidade profissional (CIP)}.
\end{itemize}

\sumario